% =====================================================================
%  Đề tài 4: Máy vector hỗ trợ và phương pháp hàm nhân
%  Biên dịch bằng XeLaTeX (xem report/README.md). Trên Overleaf:
%  Menu → Compiler → XeLaTeX.
%  File này do P1 quản lý. Mỗi người chỉ sửa file trong sections/ của mình.
% =====================================================================
\documentclass[12pt,a4paper]{article}

% ---------- Ngôn ngữ và font ----------
\usepackage{fontspec}
\usepackage{babel}
\babelprovide[import, main]{vietnamese}

% ---------- Trang ----------
\usepackage[margin=2.5cm]{geometry}
\linespread{1.15}

% ---------- Toán ----------
\usepackage{amsmath, amssymb, amsthm, mathtools, bm}

% ---------- Hình, bảng, code ----------
\usepackage{graphicx}
\graphicspath{{figures/}}
\usepackage{xcolor}
\usepackage{booktabs}
\usepackage{caption}
\usepackage{subcaption}
\usepackage{float}
\usepackage{listings}
\lstset{basicstyle=\ttfamily\small, breaklines=true, frame=single, columns=fullflexible}

% Pseudo-code: cần gói algorithms + algorithmicx (không có trong TeX Live basic).
% Cài: sudo tlmgr install algorithms algorithmicx   (Overleaf có sẵn).
\IfFileExists{algpseudocode.sty}{%
  \usepackage{algorithm}
  \usepackage{algpseudocode}
  \floatname{algorithm}{Thuật toán}
}{%
  \typeout{*** Thiếu algpseudocode.sty: chưa dùng được môi trường algorithm ***}
}

% ---------- Tài liệu tham khảo ----------
\usepackage[numbers, sort&compress]{natbib}
\renewcommand{\bibsection}{\section*{Tài liệu tham khảo}\addcontentsline{toc}{section}{Tài liệu tham khảo}}

% ---------- Liên kết (để gần cuối) ----------
\usepackage[colorlinks=true, linkcolor=blue!60!black, citecolor=green!40!black, urlcolor=blue!60!black]{hyperref}

% ---------- Định lý ----------
\theoremstyle{plain}
\newtheorem{theorem}{Định lý}[section]
\newtheorem{lemma}[theorem]{Bổ đề}
\newtheorem{proposition}[theorem]{Mệnh đề}
\newtheorem{corollary}[theorem]{Hệ quả}
\theoremstyle{definition}
\newtheorem{definition}[theorem]{Định nghĩa}
\newtheorem{example}[theorem]{Ví dụ}
\theoremstyle{remark}
\newtheorem{remark}[theorem]{Nhận xét}

% ---------- Ký hiệu dùng chung ----------
% =====================================================================
%  Ký hiệu dùng chung cho cả báo cáo (P1 quản lý).
%  Quy ước: vector chữ thường in đậm, ma trận chữ hoa in đậm,
%  vô hướng chữ thường nghiêng. Muốn thêm macro → báo P1, đừng tự định nghĩa
%  trong file section để tránh trùng tên.
% =====================================================================

% ---- Tập số, không gian ----
\newcommand{\R}{\mathbb{R}}
\newcommand{\N}{\mathbb{N}}
\newcommand{\X}{\mathcal{X}}               % không gian đầu vào
\newcommand{\Hk}{\mathcal{H}_k}            % RKHS của hàm nhân k

% ---- Vector, ma trận ----
\newcommand{\vect}[1]{\bm{#1}}
\newcommand{\vw}{\vect{w}}
\newcommand{\vx}{\vect{x}}
\newcommand{\vz}{\vect{z}}
\newcommand{\vy}{\vect{y}}
\newcommand{\vu}{\vect{u}}
\newcommand{\vc}{\vect{c}}
\newcommand{\valpha}{\vect{\alpha}}
\newcommand{\vxi}{\vect{\xi}}
\newcommand{\vmu}{\vect{\mu}}
\newcommand{\ve}{\vect{e}}                 % vector toàn số 1
\newcommand{\vzero}{\vect{0}}
\newcommand{\mK}{\mathbf{K}}               % ma trận Gram, K_ij = k(x_i, x_j)
\newcommand{\mQ}{\mathbf{Q}}               % Q_ij = y_i y_j K_ij
\newcommand{\mI}{\mathbf{I}}
\newcommand{\T}{^{\top}}                   % chuyển vị: \vw\T\vx

% ---- Chuẩn, tích vô hướng ----
\newcommand{\norm}[1]{\left\lVert #1 \right\rVert}
\newcommand{\abs}[1]{\left\lvert #1 \right\rvert}
\newcommand{\inner}[2]{\left\langle #1, #2 \right\rangle}
\newcommand{\Hnorm}[1]{\norm{#1}_{\Hk}}
\newcommand{\Hinner}[2]{\inner{#1}{#2}_{\Hk}}

% ---- Toán tử ----
\DeclareMathOperator*{\argmin}{arg\,min}
\DeclareMathOperator*{\argmax}{arg\,max}
\DeclareMathOperator{\sign}{sign}
\DeclareMathOperator{\tr}{tr}
\DeclareMathOperator{\spn}{span}
\newcommand{\ind}[1]{\mathbf{1}\!\left[#1\right]}   % hàm chỉ thị
\newcommand{\hinge}[1]{\max\{0,\, 1 - #1\}}         % mất mát hinge

% ---- SVM, SMO ----
\newcommand{\Iup}{I_{\mathrm{up}}}
\newcommand{\Ilow}{I_{\mathrm{low}}}
\newcommand{\SV}{\mathcal{S}}               % tập chỉ số vector hỗ trợ {i : α_i > 0}
\newcommand{\Dual}{D}                       % hàm mục tiêu đối ngẫu (dạng max)
\newcommand{\nsv}{n_{\mathrm{SV}}}          % số vector hỗ trợ


% ---------- Theo dõi tiến độ ----------
\newcommand{\todo}[1]{\textcolor{red}{\textbf{[TODO:} #1\textbf{]}}}
\newcommand{\owner}[1]{\marginpar{\footnotesize\textcolor{gray}{#1}}}

% =====================================================================
\title{\textbf{Đề tài 4: Máy vector hỗ trợ và phương pháp hàm nhân}\\[4pt]
       \large Bài tập lớn môn Học máy}
\author{%
  \todo{Họ tên P1} \and \todo{Họ tên P2} \and \todo{Họ tên P3}\\[6pt]
  \small Giảng viên hướng dẫn: TS. Phạm Văn Khánh}
\date{Tháng 12 năm 2026}

\begin{document}
\maketitle
\tableofcontents
\newpage

% =====================================================================
%  Giới thiệu + bảng ký hiệu (P1)
% =====================================================================
\section{Giới thiệu}\label{sec:intro}

Trong bài toán phân loại nhị phân, khi hai lớp tách được bằng một siêu phẳng thì thường có vô số
siêu phẳng như vậy. Máy vector hỗ trợ (SVM) chọn siêu phẳng có \emph{lề} lớn nhất, tức là cách xa
nhất các điểm gần nó của cả hai lớp; phiên bản lề mềm cho phép một số mẫu vi phạm lề với chi phí
điều khiển bởi tham số $C$ \citep{cortes1995support}. Hai tính chất làm nên sức mạnh của SVM.
Thứ nhất, bài toán huấn luyện là quy hoạch toàn phương lồi, nên mọi cực tiểu địa phương đều là cực
tiểu toàn cục và kết quả không phụ thuộc vào khởi tạo hay thứ tự duyệt mẫu. Thứ hai, bài toán đối ngẫu chỉ truy
cập dữ liệu qua tích vô hướng; thay tích vô hướng bằng một hàm nhân xác định dương cho phép học biên
quyết định phi tuyến trong một không gian đặc trưng, có thể vô hạn chiều, mà không cần tính ánh xạ
đặc trưng một cách tường minh \citep{scholkopf2002learning}. Nghiệm cuối cùng chỉ phụ thuộc vào một
phần nhỏ các mẫu, gọi là \emph{vector hỗ trợ}.

Những tính chất này khiến SVM đặc biệt phù hợp với bài toán có \emph{dữ liệu cỡ nhỏ} (vài trăm đến
vài nghìn mẫu) và \emph{đặc trưng thiết kế thủ công} (vài chục chiều). Đây đúng là bối cảnh của ứng
dụng trong báo cáo: phân loại giống gạo Việt Nam từ đặc trưng hình thái, màu sắc và kết cấu của
từng hạt, với mục đích hỗ trợ kiểm tra tính xác thực của nhãn các giống gạo đặc sản. Ở quy mô này,
các mô hình nhiều tham số thường cần lượng dữ liệu lớn hơn nhiều so với mức có thể thu thập thủ
công. Ngược lại, SVM chỉ có hai siêu tham số chính ($C$ và tham số của hàm nhân) chọn được bằng
kiểm định chéo, và có các kết quả lý thuyết liên hệ độ thưa của nghiệm với sai số tổng quát hóa
(Mục~\ref{sec:B6}).

Mục tiêu của báo cáo là đi trọn từ lý thuyết, qua cài đặt, tới ứng dụng:
\begin{itemize}
  \item \textbf{Lý thuyết} (Mục~\ref{sec:theory}): chứng minh đầy đủ các kết quả nền tảng, gồm hình
        học của lề cực đại, bài toán đối ngẫu và điều kiện KKT của SVM lề mềm, sự tương đương với
        cực tiểu hóa mất mát hinge, định lý biểu diễn \citep{scholkopf2001generalized}, tính chất đóng
        của hàm nhân xác định dương, và cận sai số bỏ-một theo số vector hỗ trợ.
  \item \textbf{Cài đặt} (Mục~\ref{sec:impl}): cài đặt từ đầu thuật toán SMO \citep{platt1998smo}
        với chiến lược chọn cặp vi phạm KKT lớn nhất, hạ tọa độ đối ngẫu cho SVM tuyến tính
        \citep{hsieh2008dual}, hồi quy Ridge hàm nhân, mở rộng đa lớp và Platt scaling
        \citep{platt1999probabilistic}; đối chiếu kết quả với LIBSVM \citep{chang2011libsvm}.
  \item \textbf{Bài tập và thực nghiệm} (Mục~\ref{sec:exercises}): giải tay các ví dụ nhỏ, khảo sát
        ảnh hưởng của $C$ và $\gamma$ bằng kiểm định chéo lồng nhau, kiểm chứng thực nghiệm cận sai số
        bỏ-một, và so sánh thời gian huấn luyện của SMO với hạ tọa độ đối ngẫu.
  \item \textbf{Ứng dụng} (Mục~\ref{sec:app}): thu thập ảnh hạt gạo có kiểm soát, xử lý ảnh, trích
        chọn đặc trưng, so sánh SVM với k-NN và hồi quy logistic, đánh giá theo phiên chụp để tránh
        rò rỉ dữ liệu giữa tập huấn luyện và tập kiểm tra.
\end{itemize}
Mục~\ref{sec:conclusion} tóm tắt kết quả và hướng phát triển; Phụ lục~\ref{sec:contrib} trình bày
phân công công việc. Mã nguồn (thư viện \texttt{svmlib/}) và các script tái lập thí nghiệm
(\texttt{experiments/}) nằm trong kho mã của nhóm.

\subsection{Ký hiệu}\label{sec:notation}

\begin{table}[H]
\centering
\caption{Ký hiệu dùng trong báo cáo}\label{tab:notation}
\begin{tabular}{@{}ll@{}}
\toprule
Ký hiệu & Ý nghĩa \\
\midrule
$n$, $d$ & số mẫu huấn luyện, số chiều đặc trưng \\
$\vx_i \in \R^d$, $y_i \in \{-1, +1\}$ & mẫu thứ $i$ và nhãn của nó \\
$\vw$, $b$ & vector pháp tuyến và hệ số chặn của siêu phẳng \\
$f(\vx) = \vw\T\vx + b$ & hàm quyết định; dự đoán $\sign f(\vx)$ \\
$\xi_i \ge 0$ & biến bù (slack) của mẫu $i$ trong SVM lề mềm \\
$C > 0$ & tham số phạt của SVM lề mềm \\
$\lambda > 0$ & hệ số chính quy hóa, $C = 1/(\lambda n)$ \\
$\valpha \in \R^n$ & nhân tử Lagrange (biến đối ngẫu) \\
$k(\vx, \vz)$ & hàm nhân xác định dương \\
$\mK$, $K_{ij} = k(\vx_i, \vx_j)$ & ma trận Gram \\
$\mQ$, $Q_{ij} = y_i y_j K_{ij}$ & ma trận của bài toán đối ngẫu \\
$\Hk$, $\Hnorm{\cdot}$ & RKHS sinh bởi $k$ và chuẩn trên đó \\
$\gamma$ & tham số độ rộng của hàm nhân Gauss $\exp(-\gamma\norm{\vx - \vz}^2)$ \\
$\SV = \{i : \alpha_i > 0\}$, $\nsv = \abs{\SV}$ & tập và số vector hỗ trợ \\
$\Iup$, $\Ilow$ & tập chỉ số dùng trong điều kiện KKT của SMO \\
\bottomrule
\end{tabular}
\end{table}
         % P1
% =====================================================================
%  Phần B: Chứng minh lý thuyết (P1)
%  Phát biểu và chứng minh đầy đủ của B1–B6.
% =====================================================================
\section{Cơ sở lý thuyết}\label{sec:theory}
\owner{P1}

% ---------------------------------------------------------------------
\subsection{B1. Khoảng cách tới siêu phẳng và bài toán lề cứng}\label{sec:B1}

\begin{proposition}\label{prop:B1-distance}
Cho $\vw \in \R^d \setminus \{\vzero\}$, $b \in \R$ và siêu phẳng
$H = \{\vz \in \R^d : \vw\T\vz + b = 0\}$. Với mọi $\vx \in \R^d$,
\[
  \operatorname{dist}(\vx, H) \coloneqq \min_{\vz \in H} \norm{\vx - \vz}
  = \frac{\abs{\vw\T\vx + b}}{\norm{\vw}}.
\]
\end{proposition}

\begin{proof}
Ta chứng minh hai chiều: mọi điểm của $H$ cách $\vx$ ít nhất một lượng
$\abs{\vw\T\vx + b}/\norm{\vw}$, và có một điểm của $H$ đạt đúng khoảng cách đó.

\emph{Cận dưới.} Lấy $\vz \in H$ bất kỳ. Vì $\vw\T\vz + b = 0$ nên
\[
  \vw\T\vx + b = (\vw\T\vx + b) - (\vw\T\vz + b) = \vw\T(\vx - \vz).
\]
Áp dụng bất đẳng thức Cauchy--Schwarz cho hai vector $\vw$ và $\vx - \vz$:
\[
  \abs{\vw\T\vx + b} = \abs{\vw\T(\vx - \vz)} \le \norm{\vw}\,\norm{\vx - \vz}.
\]
Vì $\vw \ne \vzero$ nên $\norm{\vw} > 0$, chia hai vế cho $\norm{\vw}$ ta được
$\norm{\vx - \vz} \ge \abs{\vw\T\vx + b}/\norm{\vw}$ với mọi $\vz \in H$.

\emph{Dấu bằng đạt được.} Đặt $t = (\vw\T\vx + b)/\norm{\vw}^2$ và $\vz^\star = \vx - t\vw$,
tức là đi từ $\vx$ dọc theo pháp tuyến $\vw$. Khi đó
\[
  \vw\T\vz^\star + b = \vw\T\vx + b - t\,\vw\T\vw
  = (\vw\T\vx + b) - \frac{\vw\T\vx + b}{\norm{\vw}^2}\,\norm{\vw}^2 = 0,
\]
nên $\vz^\star \in H$, và
\[
  \norm{\vx - \vz^\star} = \norm{t\vw} = \abs{t}\,\norm{\vw}
  = \frac{\abs{\vw\T\vx + b}}{\norm{\vw}^2}\,\norm{\vw} = \frac{\abs{\vw\T\vx + b}}{\norm{\vw}}.
\]
Hai phần trên cho thấy giá trị nhỏ nhất tồn tại, đạt tại $\vz^\star$, và bằng
$\abs{\vw\T\vx + b}/\norm{\vw}$.
\end{proof}

\begin{remark}
Đại lượng có dấu $(\vw\T\vx + b)/\norm{\vw}$ dương khi $\vx$ nằm về phía mà $\vw$ chỉ tới,
âm ở phía ngược lại. Vì vậy với nhãn $y_i \in \{-1, +1\}$, số
$y_i(\vw\T\vx_i + b)/\norm{\vw}$ bằng khoảng cách từ $\vx_i$ tới $H$ nếu $\vx_i$ được
phân loại đúng, và mang dấu âm nếu bị phân loại sai. Đây là \emph{lề hình học} của mẫu $i$.
\end{remark}

\begin{proposition}\label{prop:B1-hardmargin}
Giả sử dữ liệu $\{(\vx_i, y_i)\}_{i=1}^n$ khả tách tuyến tính và mỗi lớp có ít nhất một mẫu.
Bài toán cực đại lề hình học
\begin{equation}\label{eq:B1-maxmargin}
  \max_{\vw \ne \vzero,\, b} \; \min_{1 \le i \le n} \frac{y_i(\vw\T\vx_i + b)}{\norm{\vw}}
\end{equation}
tương đương với bài toán
\begin{equation}\label{eq:B1-hardmargin}
  \min_{\vw, b} \; \frac{1}{2}\norm{\vw}^2
  \quad \text{với ràng buộc} \quad y_i(\vw\T\vx_i + b) \ge 1, \; i = 1, \dots, n,
\end{equation}
theo nghĩa: siêu phẳng tối ưu của hai bài toán trùng nhau, và lề cực đại bằng $1/\norm{\vw^\star}$.
\end{proposition}

\begin{proof}
Với $\vw \ne \vzero$, ký hiệu lề hình học
$\gamma(\vw, b) = \min_i y_i(\vw\T\vx_i + b)/\norm{\vw}$ và tập ràng buộc của \eqref{eq:B1-hardmargin}
\[
  F = \{(\vw, b) : y_i(\vw\T\vx_i + b) \ge 1, \; i = 1, \dots, n\}.
\]
Chọn một chỉ số $p$ có $y_p = +1$ và một chỉ số $q$ có $y_q = -1$.

\emph{Bước 1 (tính chất của $F$).}
(a) $F \ne \emptyset$: vì dữ liệu khả tách nên có $(\tilde\vw, \tilde b)$ với
$\tilde s = \min_i y_i(\tilde\vw\T\vx_i + \tilde b) > 0$; khi đó
$(\tilde\vw/\tilde s, \tilde b/\tilde s) \in F$.
(b) Mọi $(\vw, b) \in F$ đều có $\vw \ne \vzero$: nếu $\vw = \vzero$ thì ràng buộc của $p$ và $q$
cho $b \ge 1$ và $-b \ge 1$, mâu thuẫn.
(c) $F$ là giao của $n$ nửa không gian đóng nên là tập lồi, đóng.

\emph{Bước 2 (bất biến tỉ lệ).} Với mọi $c > 0$,
\[
  \gamma(c\vw, cb) = \min_i \frac{c\, y_i(\vw\T\vx_i + b)}{c\norm{\vw}} = \gamma(\vw, b).
\]
Điều kiện $c > 0$ là cần thiết: với $c < 0$ mọi dấu $y_i(\cdot)$ bị đảo.
Hệ quả: nếu $\gamma(\vw, b) > 0$, đặt $s = \min_i y_i(\vw\T\vx_i + b) > 0$ thì
$(\bar\vw, \bar b) = (\vw/s, b/s) \in F$, $\min_i y_i(\bar\vw\T\vx_i + \bar b) = 1$ và
$\gamma(\vw, b) = \gamma(\bar\vw, \bar b) = 1/\norm{\bar\vw}$.

\emph{Bước 3 (\eqref{eq:B1-hardmargin} có nghiệm).} Lấy $(\vw_0, b_0) \in F$ và $R = \norm{\vw_0}$.
Với $(\vw, b) \in F$ thỏa $\norm{\vw} \le R$, ràng buộc của $p$ và $q$ cùng Cauchy--Schwarz cho
\[
  1 - R\norm{\vx_p} \;\le\; 1 - \vw\T\vx_p \;\le\; b \;\le\; -1 - \vw\T\vx_q \;\le\; -1 + R\norm{\vx_q}.
\]
Vậy $S = \{(\vw, b) \in F : \norm{\vw} \le R\}$ đóng, bị chặn, khác rỗng, tức là compact,
nên hàm liên tục $\frac12\norm{\vw}^2$ đạt giá trị nhỏ nhất trên $S$ (định lý Weierstrass).
Mọi điểm của $F \setminus S$ có $\frac12\norm{\vw}^2 > \frac12 R^2 \ge \frac12\norm{\vw_0}^2$,
nên điểm cực tiểu trên $S$ cũng là nghiệm của \eqref{eq:B1-hardmargin}.

\emph{Bước 4 (tại nghiệm, có ràng buộc chặt).} Gọi $(\vw^\star, b^\star)$ là một nghiệm của
\eqref{eq:B1-hardmargin} và $s^\star = \min_i y_i(\vw^{\star\top}\vx_i + b^\star) \ge 1$.
Giả sử $s^\star > 1$. Khi đó $(\vw^\star/s^\star, b^\star/s^\star) \in F$ và, vì $\vw^\star \ne \vzero$,
\[
  \tfrac12\norm{\vw^\star/s^\star}^2 = \tfrac{1}{(s^\star)^2}\cdot\tfrac12\norm{\vw^\star}^2
  < \tfrac12\norm{\vw^\star}^2,
\]
mâu thuẫn với tính tối ưu. Vậy $s^\star = 1$ và $\gamma(\vw^\star, b^\star) = 1/\norm{\vw^\star}$.

\emph{Bước 5 (nghiệm của \eqref{eq:B1-hardmargin} cực đại lề).} Lấy $\vw \ne \vzero$, $b$ bất kỳ.
Nếu $\gamma(\vw, b) \le 0$ thì hiển nhiên $\gamma(\vw, b) < \gamma(\vw^\star, b^\star)$.
Nếu $\gamma(\vw, b) > 0$, theo Bước 2 có $(\bar\vw, \bar b) \in F$ với
$\gamma(\vw, b) = 1/\norm{\bar\vw}$. Do $(\vw^\star, b^\star)$ tối ưu trên $F$ nên
$\norm{\bar\vw} \ge \norm{\vw^\star}$, suy ra $\gamma(\vw, b) \le 1/\norm{\vw^\star} = \gamma(\vw^\star, b^\star)$.
Vậy giá trị lớn nhất của \eqref{eq:B1-maxmargin} bằng $1/\norm{\vw^\star}$ và đạt tại $(\vw^\star, b^\star)$.

\emph{Bước 6 (nghiệm duy nhất, cùng một siêu phẳng).}
(a) $\vw^\star$ duy nhất: nếu $(\vw_1, b_1)$, $(\vw_2, b_2)$ cùng tối ưu với giá trị $p^\star$ thì
trung điểm của chúng thuộc $F$ (Bước 1c) và theo đẳng thức hình bình hành
\[
  \tfrac12\norm{\tfrac{\vw_1 + \vw_2}{2}}^2
  = \tfrac14\norm{\vw_1}^2 + \tfrac14\norm{\vw_2}^2 - \tfrac18\norm{\vw_1 - \vw_2}^2
  = p^\star - \tfrac18\norm{\vw_1 - \vw_2}^2,
\]
nhỏ hơn $p^\star$ nếu $\vw_1 \ne \vw_2$, mâu thuẫn.
(b) $b^\star$ duy nhất: nếu $(\vw^\star, b_1)$ và $(\vw^\star, b_2)$ cùng tối ưu với $b_1 < b_2$ thì
$(\vw^\star, \bar b)$, $\bar b = \frac{b_1 + b_2}{2}$, cũng tối ưu. Với $y_i = +1$:
$\vw^{\star\top}\vx_i + \bar b > \vw^{\star\top}\vx_i + b_1 \ge 1$; với $y_i = -1$:
$-(\vw^{\star\top}\vx_i + \bar b) > -(\vw^{\star\top}\vx_i + b_2) \ge 1$. Như vậy mọi ràng buộc đều
ngặt tại một nghiệm, trái với Bước 4.
(c) Nếu $(\vw, b)$ là một nghiệm của \eqref{eq:B1-maxmargin} thì $\gamma(\vw, b) = 1/\norm{\vw^\star} > 0$;
điểm chuẩn hóa $(\bar\vw, \bar b) \in F$ của nó có $\norm{\bar\vw} = 1/\gamma(\vw, b) = \norm{\vw^\star}$,
nên cũng là nghiệm của \eqref{eq:B1-hardmargin}, do đó $(\bar\vw, \bar b) = (\vw^\star, b^\star)$.
Tức là $(\vw, b) = s\,(\vw^\star, b^\star)$ với $s > 0$, mô tả cùng một siêu phẳng.
\end{proof}

\begin{remark}
Bài toán \eqref{eq:B1-maxmargin} có vô số nghiệm (cả tia $\{c(\vw^\star, b^\star) : c > 0\}$),
không lồi và không khả vi. Việc cố định lề hàm $\min_i y_i(\vw\T\vx_i + b) = 1$ chọn ra đúng
một đại diện trên tia đó và biến bài toán thành quy hoạch toàn phương lồi \eqref{eq:B1-hardmargin}.
Hệ số $\frac12$ và việc dùng $\norm{\vw}^2$ thay cho $\norm{\vw}$ chỉ để hàm mục tiêu khả vi và có
đạo hàm gọn; chúng không làm thay đổi nghiệm vì $t \mapsto \frac12 t^2$ tăng trên $[0, \infty)$.
\end{remark}

% ---------------------------------------------------------------------
\subsection{B2. Bài toán đối ngẫu của SVM lề mềm và điều kiện KKT}\label{sec:B2}

Với ánh xạ đặc trưng $\phi$ thỏa $k(\vx, \vz) = \inner{\phi(\vx)}{\phi(\vz)}$,
bài toán gốc của SVM lề mềm là
\begin{equation}\label{eq:B2-primal}
  \min_{\vw, b, \vxi} \; \frac{1}{2}\norm{\vw}^2 + C\sum_{i=1}^n \xi_i
  \quad \text{với ràng buộc} \quad
  y_i\bigl(\vw\T\phi(\vx_i) + b\bigr) \ge 1 - \xi_i, \;\; \xi_i \ge 0.
\end{equation}

\begin{theorem}\label{thm:B2-dual}
Bài toán đối ngẫu của \eqref{eq:B2-primal} là
\begin{equation}\label{eq:B2-dual}
  \max_{\valpha} \; \Dual(\valpha) = \sum_{i=1}^n \alpha_i
  - \frac{1}{2}\sum_{i,j=1}^n \alpha_i\alpha_j y_i y_j k(\vx_i, \vx_j)
  \quad \text{với ràng buộc} \quad
  0 \le \alpha_i \le C, \;\; \sum_{i=1}^n \alpha_i y_i = 0,
\end{equation}
và tại nghiệm $\vw^\star = \sum_i \alpha_i^\star y_i \phi(\vx_i)$, tức
$f(\vx) = \sum_i \alpha_i^\star y_i k(\vx_i, \vx) + b^\star$.
Hơn nữa, với mọi $i$:
\[
  \alpha_i^\star = 0 \Rightarrow y_i f(\vx_i) \ge 1, \qquad
  0 < \alpha_i^\star < C \Rightarrow y_i f(\vx_i) = 1, \qquad
  \alpha_i^\star = C \Rightarrow y_i f(\vx_i) \le 1.
\]
\end{theorem}

\begin{proof}
Để gọn, viết $\phi_i = \phi(\vx_i)$ và giả sử $\phi$ nhận giá trị trong một không gian Hilbert
(chẳng hạn $\R^D$). Bài toán \eqref{eq:B2-primal} có hàm mục tiêu lồi và mọi ràng buộc đều
affine theo $(\vw, b, \vxi)$.

\emph{Bước 1 (Lagrangian).} Gán nhân tử $\alpha_i \ge 0$ cho ràng buộc
$y_i(\vw\T\phi_i + b) \ge 1 - \xi_i$ và $\mu_i \ge 0$ cho ràng buộc $\xi_i \ge 0$:
\[
  L(\vw, b, \vxi, \valpha, \vmu) = \frac12\norm{\vw}^2 + C\sum_i \xi_i
  - \sum_i \alpha_i\bigl[y_i(\vw\T\phi_i + b) - 1 + \xi_i\bigr] - \sum_i \mu_i\xi_i .
\]
Gom các số hạng theo từng biến:
\[
  L = \Bigl(\frac12\norm{\vw}^2 - \vw\T\vu\Bigr) - b\sum_i \alpha_i y_i
  + \sum_i \xi_i (C - \alpha_i - \mu_i) + \sum_i \alpha_i,
  \qquad \vu \coloneqq \sum_i \alpha_i y_i \phi_i .
\]

\emph{Bước 2 (hàm đối ngẫu).} Hàm đối ngẫu là $g(\valpha, \vmu) = \inf_{\vw, b, \vxi} L$.
\begin{itemize}
  \item Theo $b$: $L$ tuyến tính theo $b$, nên $\inf_b L = -\infty$ trừ khi $\sum_i \alpha_i y_i = 0$.
  \item Theo $\xi_i$: tương tự, $\inf_{\xi_i} L = -\infty$ trừ khi $C - \alpha_i - \mu_i = 0$.
  \item Theo $\vw$: hoàn thành bình phương,
  $\frac12\norm{\vw}^2 - \vw\T\vu = \frac12\norm{\vw - \vu}^2 - \frac12\norm{\vu}^2$,
  nhỏ nhất tại $\vw = \vu$ với giá trị $-\frac12\norm{\vu}^2$.
\end{itemize}
Vậy khi $\sum_i \alpha_i y_i = 0$ và $\mu_i = C - \alpha_i$,
\[
  g(\valpha, \vmu) = \sum_i \alpha_i - \frac12\Bigl\lVert \sum_i \alpha_i y_i \phi_i \Bigr\rVert^2
  = \sum_i \alpha_i - \frac12\sum_{i,j} \alpha_i \alpha_j y_i y_j k(\vx_i, \vx_j) = \Dual(\valpha),
\]
và $g = -\infty$ trong mọi trường hợp khác. Điều kiện $\mu_i = C - \alpha_i \ge 0$ cùng
$\alpha_i \ge 0$ tương đương $0 \le \alpha_i \le C$. Cực đại $g$ chính là bài toán \eqref{eq:B2-dual}.

\emph{Bước 3 (đối ngẫu mạnh và KKT).} Bài toán gốc lồi, ràng buộc affine và khả thi
(ví dụ $\vw = \vzero$, $b = 0$, $\xi_i = 1$), nên theo điều kiện Slater dạng ràng buộc affine
\citep[Mục 5.2.3]{boyd2004convex}, đối ngẫu mạnh xảy ra: giá trị tối ưu của hai bài toán bằng nhau.
Khi đó $(\vw^\star, b^\star, \vxi^\star)$ và $(\valpha^\star, \vmu^\star)$ lần lượt là nghiệm
của bài toán gốc và đối ngẫu khi và chỉ khi chúng thỏa các điều kiện KKT
\citep[Mục 5.5.3]{boyd2004convex}:
\begin{enumerate}
  \item khả thi gốc: $y_i f(\vx_i) \ge 1 - \xi_i^\star$ và $\xi_i^\star \ge 0$,
        với $f(\vx) = \vw^{\star\top}\phi(\vx) + b^\star$;
  \item khả thi đối ngẫu: $\alpha_i^\star \ge 0$, $\mu_i^\star \ge 0$;
  \item điều kiện dừng: $\vw^\star = \sum_i \alpha_i^\star y_i \phi_i$, $\sum_i \alpha_i^\star y_i = 0$,
        $\alpha_i^\star + \mu_i^\star = C$;
  \item bù lệch: $\alpha_i^\star\bigl[y_i f(\vx_i) - 1 + \xi_i^\star\bigr] = 0$ và $\mu_i^\star \xi_i^\star = 0$.
\end{enumerate}
Từ điều kiện dừng,
$f(\vx) = \sum_i \alpha_i^\star y_i \phi_i\T\phi(\vx) + b^\star = \sum_i \alpha_i^\star y_i k(\vx_i, \vx) + b^\star$.

\emph{Bước 4 (ba trường hợp).}
\begin{itemize}
  \item $\alpha_i^\star = 0$: khi đó $\mu_i^\star = C > 0$, bù lệch cho $\xi_i^\star = 0$,
        và khả thi gốc cho $y_i f(\vx_i) \ge 1$.
  \item $0 < \alpha_i^\star < C$: khi đó $\mu_i^\star = C - \alpha_i^\star > 0$ nên $\xi_i^\star = 0$;
        vì $\alpha_i^\star > 0$, bù lệch cho $y_i f(\vx_i) = 1 - \xi_i^\star = 1$.
  \item $\alpha_i^\star = C$: vì $\alpha_i^\star > 0$, bù lệch cho
        $y_i f(\vx_i) = 1 - \xi_i^\star \le 1$ do $\xi_i^\star \ge 0$. \qedhere
\end{itemize}
\end{proof}

\begin{remark}\label{rem:B2-kkt}
Chiều ngược lại cũng suy ra được từ KKT: nếu $y_i f(\vx_i) > 1$ thì
$y_i f(\vx_i) - 1 + \xi_i^\star > 0$ nên $\alpha_i^\star = 0$; nếu $y_i f(\vx_i) < 1$ thì
$\xi_i^\star \ge 1 - y_i f(\vx_i) > 0$ nên $\mu_i^\star = 0$, tức $\alpha_i^\star = C$.
Như vậy mẫu nằm ngoài lề không ảnh hưởng tới nghiệm; chỉ các \emph{vector hỗ trợ}
($\alpha_i^\star > 0$) xuất hiện trong $f$. Hệ số chặn được tính từ một vector hỗ trợ tự do
bất kỳ ($0 < \alpha_i^\star < C$): $b^\star = y_i - \sum_j \alpha_j^\star y_j k(\vx_j, \vx_i)$.
\end{remark}

% ---------------------------------------------------------------------
\subsection{B3. Dạng hàm mất mát hinge}\label{sec:B3}

\begin{proposition}\label{prop:B3-hinge}
Với $\lambda = 1/(Cn)$, $(\vw^\star, b^\star)$ là nghiệm của \eqref{eq:B2-primal}
khi và chỉ khi nó là nghiệm của
\begin{equation}\label{eq:B3-erm}
  \min_{\vw, b} \; \frac{\lambda}{2}\norm{\vw}^2
  + \frac{1}{n}\sum_{i=1}^n \hinge{y_i f(\vx_i)}.
\end{equation}
Ngoài ra, hàm $t \mapsto \hinge{t}$ lồi và $\hinge{t} \ge \ind{t \le 0}$ với mọi $t \in \R$.
\end{proposition}

\begin{proof}
Ký hiệu $h(t) = \hinge{t}$ và $f(\vx) = \vw\T\phi(\vx) + b$.

\emph{Bước 1 (khử biến bù).} Cố định $(\vw, b)$. Hai ràng buộc của $\xi_i$ trong
\eqref{eq:B2-primal} là $\xi_i \ge 1 - y_i f(\vx_i)$ và $\xi_i \ge 0$, tương đương
$\xi_i \ge h(y_i f(\vx_i))$. Hàm mục tiêu tăng ngặt theo mỗi $\xi_i$ (vì $C > 0$), nên với
$(\vw, b)$ cố định, $\vxi$ tốt nhất là $\xi_i = h(y_i f(\vx_i))$. Do đó $(\vw^\star, b^\star, \vxi^\star)$
là nghiệm của \eqref{eq:B2-primal} khi và chỉ khi $\xi_i^\star = h(y_i f^\star(\vx_i))$ và
$(\vw^\star, b^\star)$ cực tiểu hóa
\[
  P(\vw, b) = \frac12\norm{\vw}^2 + C\sum_{i=1}^n h\bigl(y_i f(\vx_i)\bigr).
\]

\emph{Bước 2 (đổi tỉ lệ).} Với $\lambda = 1/(Cn)$,
\[
  P(\vw, b) = Cn\Bigl[\frac{1}{Cn}\cdot\frac12\norm{\vw}^2 + \frac1n\sum_{i=1}^n h\bigl(y_i f(\vx_i)\bigr)\Bigr]
  = Cn \cdot J_\lambda(\vw, b),
\]
trong đó $J_\lambda$ là hàm mục tiêu của \eqref{eq:B3-erm}. Nhân với hằng số dương $Cn$
không đổi tập nghiệm, nên $P$ và $J_\lambda$ có cùng nghiệm.

\emph{Bước 3 (tính lồi).} $h(t) = \max\{\ell_1(t), \ell_2(t)\}$ với $\ell_1(t) = 0$,
$\ell_2(t) = 1 - t$ là các hàm affine. Với $t, s \in \R$ và $\theta \in [0, 1]$, mỗi $\ell_k$ thỏa
$\ell_k(\theta t + (1-\theta)s) = \theta\ell_k(t) + (1-\theta)\ell_k(s) \le \theta h(t) + (1-\theta)h(s)$;
lấy max theo $k$ ta được $h(\theta t + (1-\theta)s) \le \theta h(t) + (1-\theta)h(s)$.

\emph{Bước 4 (cận trên của mất mát 0--1).} Nếu $t \le 0$ thì $h(t) \ge 1 - t \ge 1 = \ind{t \le 0}$.
Nếu $t > 0$ thì $h(t) \ge 0 = \ind{t \le 0}$.
\end{proof}

\begin{remark}
Vì $y_i f(\vx_i)$ là hàm affine của $(\vw, b)$ và $h$ lồi, $J_\lambda$ là tổng của một hàm lồi
ngặt theo $\vw$ và các hàm lồi, nên \eqref{eq:B3-erm} là bài toán lồi. Bước 4 cho
$\frac1n\sum_i \ind{y_i f(\vx_i) \le 0} \le \frac1n\sum_i h(y_i f(\vx_i))$: sai số huấn luyện bị chặn
trên bởi rủi ro hinge. Thay vì cực tiểu mất mát 0--1 (không lồi, NP-khó), SVM cực tiểu một cận
trên lồi của nó cộng với số hạng chính quy hóa $\frac\lambda2\norm{\vw}^2$. Quan hệ $C = 1/(\lambda n)$
cho thấy $C$ lớn tương ứng với chính quy hóa yếu.
\end{remark}

% ---------------------------------------------------------------------
\subsection{B4. Định lý biểu diễn}\label{sec:B4}

\begin{theorem}[\citet{scholkopf2001generalized}]\label{thm:B4-representer}
Cho $k$ là hàm nhân xác định dương trên $\X$ với RKHS $\Hk$, dữ liệu
$(\vx_1, y_1), \dots, (\vx_n, y_n)$, hàm mất mát bất kỳ
$c : (\X \times \R \times \R)^n \to \R \cup \{\infty\}$ và
$\Omega : [0, \infty) \to \R$ tăng nghiêm ngặt. Xét bài toán
\begin{equation}\label{eq:B4-risk}
  \min_{f \in \Hk} \; R(f) \coloneqq c\bigl((\vx_1, y_1, f(\vx_1)), \dots, (\vx_n, y_n, f(\vx_n))\bigr)
  + \Omega\bigl(\Hnorm{f}\bigr)
\end{equation}
và giả sử giá trị nhỏ nhất của nó hữu hạn. Khi đó mọi nghiệm $f^\star$ của \eqref{eq:B4-risk}
đều có dạng $f^\star = \sum_{i=1}^n \alpha_i k(\vx_i, \cdot)$ với $\alpha_i \in \R$.
\end{theorem}

\begin{proof}
Nhắc lại hai tính chất định nghĩa RKHS \citep{aronszajn1950theory}: với mọi $\vx \in \X$,
hàm $k(\vx, \cdot)$ thuộc $\Hk$, và với mọi $f \in \Hk$ có \emph{tính chất tái sinh}
$f(\vx) = \Hinner{f}{k(\vx, \cdot)}$.

\emph{Bước 1 (phân tích trực giao).} Đặt $V = \spn\{k(\vx_1, \cdot), \dots, k(\vx_n, \cdot)\}$.
$V$ hữu hạn chiều nên là không gian con đóng của $\Hk$, và theo định lý chiếu trực giao mọi
$f \in \Hk$ viết được duy nhất dưới dạng $f = f_\parallel + f_\perp$ với $f_\parallel \in V$ và
$f_\perp \perp V$.

\emph{Bước 2 (giá trị tại dữ liệu chỉ phụ thuộc $f_\parallel$).} Với mỗi $j$, vì
$k(\vx_j, \cdot) \in V$ nên $\Hinner{f_\perp}{k(\vx_j, \cdot)} = 0$, do đó
\[
  f(\vx_j) = \Hinner{f_\parallel + f_\perp}{k(\vx_j, \cdot)}
  = \Hinner{f_\parallel}{k(\vx_j, \cdot)} = f_\parallel(\vx_j).
\]
Vậy số hạng mất mát $c(\cdots)$ của $f$ và của $f_\parallel$ bằng nhau.

\emph{Bước 3 (chuẩn).} Theo định lý Pythagoras,
$\Hnorm{f}^2 = \Hnorm{f_\parallel}^2 + \Hnorm{f_\perp}^2 \ge \Hnorm{f_\parallel}^2$, dấu bằng
khi và chỉ khi $f_\perp = 0$.

\emph{Bước 4 (kết luận).} Giả sử $f^\star$ là một nghiệm nhưng $f^\star_\perp \ne 0$. Khi đó
$\Hnorm{f^\star_\parallel} < \Hnorm{f^\star}$ và vì $\Omega$ tăng nghiêm ngặt,
$\Omega(\Hnorm{f^\star_\parallel}) < \Omega(\Hnorm{f^\star})$. Kết hợp Bước 2 và giá trị
$R(f^\star)$ hữu hạn, ta được $R(f^\star_\parallel) < R(f^\star)$, mâu thuẫn. Vậy
$f^\star = f^\star_\parallel \in V$, tức $f^\star = \sum_i \alpha_i k(\vx_i, \cdot)$.
\end{proof}

\begin{remark}
SVM lề mềm là trường hợp riêng với $\Omega(t) = \frac12 t^2$ và
$c = C\sum_i \hinge{y_i(f(\vx_i) + b)}$. Hệ số chặn $b$ không bị chính quy hóa; phiên bản
\emph{bán tham số} của định lý \citep{scholkopf2001generalized} cho nghiệm dạng
$f^\star + b$ với $f^\star \in V$. Đây là lý do nghiệm SVM chỉ cần $n$ hệ số $\alpha_i$ và chỉ
truy cập dữ liệu qua $k(\vx_i, \vx_j)$, dù $\Hk$ có thể vô hạn chiều (như với hàm nhân Gauss).
Nếu $\Omega$ chỉ không giảm thì lập luận trên vẫn cho thấy \emph{có ít nhất một} nghiệm dạng
này (thay $f^\star$ bởi $f^\star_\parallel$), nhưng không phải mọi nghiệm.
\end{remark}

% ---------------------------------------------------------------------
\subsection{B5. Tính chất đóng của hàm nhân xác định dương}\label{sec:B5}

\begin{definition}\label{def:B5-pd}
Hàm đối xứng $k : \X \times \X \to \R$ được gọi là \emph{xác định dương} nếu với mọi
$m \in \N$, mọi $\vx_1, \dots, \vx_m \in \X$ và mọi $\vc \in \R^m$:
$\sum_{i,j=1}^m c_i c_j k(\vx_i, \vx_j) \ge 0$, tức là mọi ma trận Gram đều nửa xác định dương.
\end{definition}

\begin{lemma}[Định lý tích Schur]\label{lem:B5-schur}
Nếu $\mathbf{A}, \mathbf{B} \in \R^{m \times m}$ nửa xác định dương thì tích Hadamard
$\mathbf{A} \circ \mathbf{B}$ cũng nửa xác định dương.
\end{lemma}

\begin{proof}
Vì $\mathbf{A}$, $\mathbf{B}$ đối xứng nửa xác định dương, theo định lý phổ
$\mathbf{A} = \sum_k \lambda_k \vu_k\vu_k\T$ và $\mathbf{B} = \sum_l \mu_l \vect{v}_l\vect{v}_l\T$ với
$\lambda_k, \mu_l \ge 0$. Phần tử $(i, j)$ của $(\vu\vu\T) \circ (\vect{v}\vect{v}\T)$ là
$u_i u_j v_i v_j = (\vu \circ \vect{v})_i (\vu \circ \vect{v})_j$, nên
$(\vu\vu\T) \circ (\vect{v}\vect{v}\T) = (\vu \circ \vect{v})(\vu \circ \vect{v})\T$. Tích Hadamard
tuyến tính theo từng thừa số, do đó với mọi $\vc \in \R^m$,
\[
  \vc\T(\mathbf{A} \circ \mathbf{B})\vc
  = \sum_{k,l} \lambda_k \mu_l\, \vc\T(\vu_k \circ \vect{v}_l)(\vu_k \circ \vect{v}_l)\T\vc
  = \sum_{k,l} \lambda_k \mu_l \bigl((\vu_k \circ \vect{v}_l)\T\vc\bigr)^2 \ge 0.
\]
Hiển nhiên $\mathbf{A} \circ \mathbf{B}$ đối xứng. Vậy $\mathbf{A} \circ \mathbf{B}$ nửa xác định dương.
\end{proof}

\begin{proposition}\label{prop:B5-closure}
Cho $k_1, k_2$ là các hàm nhân xác định dương trên $\X$. Khi đó các hàm sau cũng
xác định dương: (i) $k_1 + k_2$; (ii) $c\,k_1$ với $c \ge 0$; (iii) $k_1 k_2$;
(iv) $k = \lim_{m\to\infty} k_m$ nếu mỗi $k_m$ xác định dương và giới hạn tồn tại tại mọi điểm;
(v) $\exp(k_1)$.
\end{proposition}

\begin{proof}
Cố định $m$, các điểm $\vx_1, \dots, \vx_m$ và $\vc \in \R^m$; gọi $\mK_1$, $\mK_2$ là ma trận Gram
của $k_1$, $k_2$ trên các điểm này. Tính đối xứng của các hàm mới là hiển nhiên; ta chỉ cần
kiểm tra $\vc\T\mK\vc \ge 0$ cho ma trận Gram $\mK$ tương ứng.
\begin{enumerate}
  \item[(i)] Ma trận Gram của $k_1 + k_2$ là $\mK_1 + \mK_2$, và
        $\vc\T(\mK_1 + \mK_2)\vc = \vc\T\mK_1\vc + \vc\T\mK_2\vc \ge 0$.
  \item[(ii)] $\vc\T(c\mK_1)\vc = c\,\vc\T\mK_1\vc \ge 0$ vì $c \ge 0$.
  \item[(iii)] Ma trận Gram của $k_1 k_2$ có phần tử $k_1(\vx_i, \vx_j)k_2(\vx_i, \vx_j)$, tức là
        $\mK_1 \circ \mK_2$, nửa xác định dương theo Bổ đề~\ref{lem:B5-schur}.
  \item[(iv)] Với mỗi $m'$, $\sum_{i,j} c_i c_j k_{m'}(\vx_i, \vx_j) \ge 0$. Đây là tổng hữu hạn
        của các dãy hội tụ, nên khi $m' \to \infty$ nó hội tụ về $\sum_{i,j} c_i c_j k(\vx_i, \vx_j)$.
        Giới hạn của một dãy số không âm thì không âm.
  \item[(v)] Hàm hằng $1$ xác định dương vì ma trận Gram của nó là $\ve\ve\T$ và
        $\vc\T\ve\ve\T\vc = (\sum_i c_i)^2 \ge 0$. Theo (iii) và quy nạp, $k_1^p$ xác định dương
        với mọi $p \in \N$. Theo (i), (ii), tổng riêng $s_N = \sum_{p=0}^N k_1^p / p!$ xác định dương.
        Chuỗi Taylor của hàm mũ hội tụ tại mọi điểm nên $s_N(\vx, \vz) \to \exp(k_1(\vx, \vz))$,
        và theo (iv), $\exp(k_1)$ xác định dương. \qedhere
\end{enumerate}
\end{proof}

\begin{corollary}\label{cor:B5-gauss}
Với mọi $\gamma > 0$, hàm nhân Gauss $k(\vx, \vz) = \exp(-\gamma\norm{\vx - \vz}^2)$ xác định dương.
\end{corollary}

\begin{proof}
Khai triển $\norm{\vx - \vz}^2 = \norm{\vx}^2 + \norm{\vz}^2 - 2\vx\T\vz$ được
\[
  k(\vx, \vz) = g(\vx)\, g(\vz)\, \exp(2\gamma\,\vx\T\vz), \qquad g(\vx) = \exp(-\gamma\norm{\vx}^2) > 0.
\]
Ta chứng minh từng thừa số.
\begin{enumerate}
  \item Hàm nhân tuyến tính $\vx\T\vz$ xác định dương: với ma trận dữ liệu $\mathbf{X}$ (các hàng là
        $\vx_i\T$), ma trận Gram là $\mathbf{X}\mathbf{X}\T$ và
        $\vc\T\mathbf{X}\mathbf{X}\T\vc = \norm{\mathbf{X}\T\vc}^2 \ge 0$.
  \item Theo Mệnh đề~\ref{prop:B5-closure}(ii), $2\gamma\,\vx\T\vz$ xác định dương; theo (v),
        $k'(\vx, \vz) = \exp(2\gamma\,\vx\T\vz)$ xác định dương.
  \item Hàm $g(\vx)g(\vz)$ xác định dương vì
        $\sum_{i,j} c_i c_j g(\vx_i) g(\vx_j) = \bigl(\sum_i c_i g(\vx_i)\bigr)^2 \ge 0$.
        Theo (iii), tích $g(\vx)g(\vz)\,k'(\vx, \vz) = k(\vx, \vz)$ xác định dương. \qedhere
\end{enumerate}
\end{proof}

\begin{remark}
Hàm nhân đa thức $(\vx\T\vz + c)^p$ với $c \ge 0$, $p \in \N$ cũng xác định dương theo cùng cách:
$\vx\T\vz + c$ là tổng của hàm nhân tuyến tính và hằng số không âm, rồi áp dụng (iii) $p$ lần.
Như vậy cả ba hàm nhân dùng trong Phần A đều hợp lệ.
\end{remark}

% ---------------------------------------------------------------------
\subsection{B6. Cận sai số kiểm định chéo bỏ-một}\label{sec:B6}

\begin{theorem}[\citet{vapnik1995nature}]\label{thm:B6-loo}
Giả sử dữ liệu khả tách và $\valpha$ là nghiệm đối ngẫu của SVM lề cứng huấn luyện trên
cả $n$ mẫu. Gọi $f^{(-i)}$ là SVM lề cứng huấn luyện khi bỏ mẫu $i$. Khi đó
\[
  \frac{1}{n}\sum_{i=1}^n \ind{y_i f^{(-i)}(\vx_i) \le 0}
  \;\le\; \frac{\abs{\{i : \alpha_i > 0\}}}{n}.
\]
\end{theorem}

\begin{proof}
SVM lề cứng là bài toán \eqref{eq:B1-hardmargin} với $\vx_i$ thay bởi $\phi(\vx_i)$; giả sử mỗi lớp
có ít nhất một mẫu. Lập luận ở Mục~\ref{sec:B2} (bỏ biến bù, tức ràng buộc $\alpha_i \le C$ biến mất)
cho thấy $(\vw^\star, b^\star)$ là nghiệm khi và chỉ khi tồn tại $\valpha$ thỏa KKT:
\[
  \vw^\star = \sum_j \alpha_j y_j \phi(\vx_j), \quad \sum_j \alpha_j y_j = 0, \quad \alpha_j \ge 0, \quad
  y_j f(\vx_j) \ge 1, \quad \alpha_j\bigl(y_j f(\vx_j) - 1\bigr) = 0 .
\]
Theo Mệnh đề~\ref{prop:B1-hardmargin} (Bước 6, lập luận không đổi trong không gian đặc trưng),
nghiệm $(\vw^\star, b^\star)$ là duy nhất. Đặt $\SV = \{j : \alpha_j > 0\}$.

\emph{Bước 1 ($\SV$ chứa cả hai lớp).} $\vw^\star \ne \vzero$ (Mệnh đề~\ref{prop:B1-hardmargin}, Bước 1b)
nên $\SV \ne \emptyset$. Vì $\sum_{j \in \SV} \alpha_j y_j = 0$ với mọi $\alpha_j > 0$, các $y_j$
trong $\SV$ không thể cùng dấu.

\emph{Bước 2 (bỏ một mẫu không phải vector hỗ trợ).} Lấy $i \notin \SV$, tức $\alpha_i = 0$. Tập dữ liệu
$D^{(-i)}$ (bỏ mẫu $i$) vẫn khả tách và chứa $\SV$, nên chứa cả hai lớp. Xét $(\vw^\star, b^\star)$
cùng các nhân tử $(\alpha_j)_{j \ne i}$ cho bài toán lề cứng trên $D^{(-i)}$:
\begin{itemize}
  \item $\vw^\star = \sum_{j \ne i} \alpha_j y_j \phi(\vx_j)$ và $\sum_{j \ne i}\alpha_j y_j = 0$,
        vì số hạng thứ $i$ bằng $0$;
  \item khả thi gốc, $\alpha_j \ge 0$ và bù lệch với $j \ne i$ giữ nguyên từ bài toán đầy đủ.
\end{itemize}
Bài toán lồi với ràng buộc affine nên KKT là điều kiện đủ; do đó $(\vw^\star, b^\star)$ là nghiệm, và
theo tính duy nhất là \emph{nghiệm duy nhất} trên $D^{(-i)}$. Vậy $f^{(-i)} = f$.

\emph{Bước 3 (không có lỗi tại mẫu bị bỏ).} Từ khả thi gốc trên bài toán đầy đủ,
$y_i f^{(-i)}(\vx_i) = y_i f(\vx_i) \ge 1 > 0$, nên mẫu $i$ được dự đoán đúng.

\emph{Bước 4 (đếm).} Lỗi bỏ-một chỉ có thể xảy ra tại $i \in \SV$, do đó
$\sum_i \ind{y_i f^{(-i)}(\vx_i) \le 0} \le \abs{\SV}$. Chia hai vế cho $n$.
\end{proof}

\begin{remark}\label{rem:B6-soft}
(1) Nếu nghiệm đối ngẫu không duy nhất, cận đúng với mọi nghiệm $\valpha$, nên có thể chọn nghiệm
có ít vector hỗ trợ nhất. (2) Với SVM lề mềm, lập luận giống hệt: mẫu có $\alpha_i = 0$ thỏa
$y_i f(\vx_i) \ge 1$ (Định lý~\ref{thm:B2-dual}), và bỏ nó đi thì $(\alpha_j)_{j \ne i}$ vẫn thỏa KKT
của bài toán mới. Khi nghiệm đối ngẫu có ít nhất một vector hỗ trợ tự do, $b^\star$ được xác định
duy nhất nên $f^{(-i)} = f$ và cận vẫn đúng; Mục~\ref{sec:C4} kiểm chứng điều này bằng thực nghiệm.
(3) Sai số bỏ-một là ước lượng gần không chệch của sai số kiểm tra khi huấn luyện trên $n - 1$ mẫu,
nên tỉ lệ vector hỗ trợ cho một chặn trên dễ tính của sai số tổng quát hóa \citep{vapnik1995nature}:
mô hình càng thưa (ít vector hỗ trợ) thì cận càng chặt.
\end{remark}
         % P1
% =====================================================================
%  Phần A: Phát triển phương pháp và cài đặt (P2)
%  Công thức SMO: xem notes/kkt_cheatsheet.md. Dẫn chiếu lý thuyết:
%  Định lý~\ref{thm:B2-dual} (đối ngẫu + KKT).
% =====================================================================
\section{Cài đặt thuật toán}\label{sec:impl}
\owner{P2}

\subsection{A1. Thuật toán SMO}\label{sec:A1}
\todo{Bài toán đối ngẫu (dẫn Định lý~\ref{thm:B2-dual}); chọn cặp vi phạm KKT lớn nhất
\citep{keerthi2001improvements}; cập nhật giải tích 2 biến \citep{platt1998smo};
cập nhật gradient; tính $b$; 3 hàm nhân; pseudo-code; hình hội tụ.}

\subsection{A2. Hạ tọa độ đối ngẫu và hồi quy Ridge hàm nhân}\label{sec:A2}
\todo{DCD cho SVM tuyến tính \citep{hsieh2008dual}; kernel ridge regression với
nghiệm dạng đóng $\valpha = (\mK + \lambda\mI)^{-1}\vy$.}

\subsection{A3. Mở rộng đa lớp và Platt scaling}\label{sec:A3}
\todo{Một-đấu-một, một-đấu-phần còn lại; Platt scaling \citep{platt1999probabilistic, lin2007note};
pairwise coupling \citep{wu2004probability}; reliability diagram, Brier score.}

\subsection{A4. Đối chiếu với LIBSVM}\label{sec:A4}
\todo{Bảng so sánh với LIBSVM \citep{chang2011libsvm} qua scikit-learn \citep{pedregosa2011scikit}:
hàm mục tiêu đối ngẫu, tập vector hỗ trợ, $b$, dự đoán; nêu ngưỡng sai số cho phép.}
 % P2
% =====================================================================
%  Phần C: Bài tập (P1: C1–C4, P2: C5)
% =====================================================================
\section{Bài tập và thực nghiệm}\label{sec:exercises}

% ---------------------------------------------------------------------
\subsection{C1. Giải tay SVM lề cứng}\label{sec:C1}
\owner{P1}
\textbf{Đề bài.} Lớp dương gồm $(2, 0)$ và $(0, 2)$, lớp âm gồm $(0, 0)$. Giải tay SVM lề cứng:
xác định $\vw$, $b$, độ rộng lề, các nhân tử Lagrange và kiểm tra toàn bộ điều kiện KKT.

\todo{C1: lời giải; đối chiếu với SMO của P2.}

% ---------------------------------------------------------------------
\subsection{C2. Bài toán XOR và ánh xạ đặc trưng}\label{sec:C2}
\owner{P1}
\textbf{Đề bài.} Chứng minh bài toán XOR không khả tách tuyến tính trong $\R^2$; tìm ánh xạ
đặc trưng tường minh $\phi$ với $\phi(\vx)\T\phi(\vz) = (\vx\T\vz + 1)^2$ và một siêu phẳng tách
trong không gian đặc trưng.

\todo{C2: lời giải.}

% ---------------------------------------------------------------------
\subsection{C3. Một hàm không phải hàm nhân xác định dương}\label{sec:C3}
\owner{P1}
\textbf{Đề bài.} Chứng minh $k(\vx, \vz) = \norm{\vx - \vz}$ không là hàm nhân xác định dương
bằng một phản ví dụ.

\todo{C3: lời giải.}

% ---------------------------------------------------------------------
\subsection{\texorpdfstring{C4. Ảnh hưởng của $C$ và $\gamma$}{C4. Ảnh hưởng của C và γ}}\label{sec:C4}
\owner{P1}
\todo{C4: thiết lập thí nghiệm; heatmap sai số kiểm định chéo lồng nhau;
số vector hỗ trợ và độ rộng lề theo $C$, $\gamma$; đối chiếu với cận ở Định lý~\ref{thm:B6-loo}.}

% ---------------------------------------------------------------------
\subsection{C5. Thời gian huấn luyện: SMO và hạ tọa độ đối ngẫu}\label{sec:C5}
\owner{P2}
\todo{C5: $n$ từ $10^3$ đến $10^5$, hàm nhân tuyến tính; đồ thị log-log; bảng thời gian.}
      % P1 (C1–C4), P2 (C5)
% =====================================================================
%  Phần D: Ứng dụng phân loại giống gạo Việt Nam từ ảnh hạt (P3)
% =====================================================================
\section{Ứng dụng: phân loại giống gạo Việt Nam từ ảnh hạt}\label{sec:app}
\owner{P3}

\subsection{D1. Thiết kế thu thập dữ liệu}\label{sec:D1}
\todo{Các giống gạo và nguồn gốc; quy trình chụp (giá cố định, nền, ánh sáng, thước);
số hạt mỗi giống, số phiên, số thiết bị; ảnh bộ setup.}

\subsection{D2. Xử lý ảnh}\label{sec:D2}
\todo{Phân ngưỡng Otsu \citep{otsu1979threshold}; hình thái học; watershed
\citep{vincent1991watersheds}; gán nhãn thành phần liên thông; loại hạt gãy; tỉ lệ tách đúng.}

\subsection{D3. Trích chọn đặc trưng}\label{sec:D3}
\todo{Hình thái, màu sắc (RGB, HSV, CIELAB), kết cấu GLCM \citep{haralick1973textural}; EDA.}

\subsection{D4. Mô hình hóa}\label{sec:D4}
\todo{SVM (Mục~\ref{sec:impl}) với 3 hàm nhân, k-NN, hồi quy logistic;
chuẩn hóa đặc trưng; kiểm định chéo lồng nhau theo phiên chụp.}

\subsection{D5. Đánh giá}\label{sec:D5}
\todo{Chia theo phiên; ma trận nhầm lẫn; độ bền khi đổi thiết bị; ablation;
thảo luận ứng dụng kiểm tra nhãn gạo đặc sản.}
    % P3
% =====================================================================
%  Kết luận (P1)
%  Đã viết: Phần B, C1–C4. Chờ số liệu cuối của P2 (A, C5) và P3 (D)
%  để điền các chỗ \todo.
% =====================================================================
\section{Kết luận}\label{sec:conclusion}

Báo cáo đã đi trọn chuỗi từ cơ sở lý thuyết của SVM và phương pháp hàm nhân, qua cài đặt từ đầu,
tới thực nghiệm và ứng dụng. Các kết quả chính theo từng phần như sau.

\paragraph{Lý thuyết (Phần B).} Mục~\ref{sec:theory} chứng minh đầy đủ sáu kết quả nền tảng. Bài toán
cực đại lề tương đương quy hoạch toàn phương lồi $\min \frac12\norm{\vw}^2$ và có nghiệm duy nhất; bài
toán đối ngẫu của SVM lề mềm chỉ phụ thuộc dữ liệu qua hàm nhân, và điều kiện KKT phân loại mỗi mẫu
thành ngoài lề, trên lề hoặc vi phạm lề; SVM lề mềm chính là cực tiểu hóa rủi ro hinge có chính quy
hóa với $C = 1/(\lambda n)$; định lý biểu diễn giải thích vì sao nghiệm chỉ cần $n$ hệ số; tổng, tích,
giới hạn và hàm mũ của hàm nhân xác định dương vẫn xác định dương, kéo theo tính hợp lệ của hàm nhân
Gauss; và tỉ lệ vector hỗ trợ chặn trên sai số bỏ-một. Trong quá trình chứng minh, hai giả thiết ngầm
của đề bài được nêu rõ: mỗi lớp có ít nhất một mẫu (B1, B6) và giá trị tối ưu hữu hạn (B4).

\paragraph{Cài đặt (Phần A).} \todo{P2: tóm tắt 2--3 câu kết quả của A1--A4: SMO, DCD, KRR, đa lớp,
Platt scaling; mức sai lệch so với LIBSVM (hàm mục tiêu đối ngẫu, tập vector hỗ trợ, dự đoán).}

\paragraph{Bài tập và thực nghiệm (Phần C).} Lời giải tay của C1 ($\vw = (1, 1)$, $b = -1$,
$\valpha = (\frac12, \frac12, 1)$, độ rộng lề $\sqrt2$) và C2 (siêu phẳng tách XOR trong không gian
đặc trưng của hàm nhân đa thức bậc hai, trùng với nghiệm lề cực đại $f(\vx) = -x_1x_2$) được SMO của
nhóm và LIBSVM tìm lại chính xác. Khảo sát $C$ và $\gamma$ (C4) bằng kiểm định chéo lồng nhau cho sai
số kiểm tra $7{,}0 \pm 1{,}1\%$ trên moons và $1{,}9 \pm 1{,}0\%$ trên breast cancer. Thực nghiệm xác
nhận các dự đoán của lý thuyết: khi $C$ tăng, độ rộng lề và số vector hỗ trợ đều không tăng với mọi
$\gamma$; $\gamma$ quá lớn dẫn tới ma trận Gram gần đơn vị và mô hình chỉ ghi nhớ dữ liệu huấn luyện;
và cận $\text{LOO} \le \nsv/n$ đúng ở cả $62$ cấu hình đã thử. \todo{P2: 1--2 câu kết quả C5
(thời gian SMO so với DCD khi $n$ tăng).}

\paragraph{Ứng dụng (Phần D).} \todo{P3: tóm tắt 2--3 câu: dữ liệu thu thập được, mô hình tốt nhất
và sai số khi chia theo phiên chụp, độ bền khi đổi thiết bị, khả năng dùng để kiểm tra nhãn gạo đặc sản.}

\paragraph{Hạn chế.}
\begin{itemize}
  \item Cận bỏ-một của B6 thường lỏng: trên các cấu hình đã thử, sai số bỏ-một thực tế chỉ bằng khoảng
        $0{,}36$--$0{,}37$ lần cận (trung vị). Với SVM lề mềm, cận chỉ được bảo đảm khi nghiệm có vector
        hỗ trợ tự do; khi mọi $\alpha_i = C$, sai số bỏ-một có thể lên tới $100\%$ dù sai số kiểm định
        chéo chỉ khoảng $17\%$.
  \item Khảo sát C4 dùng hai bộ dữ liệu chuẩn và lưới thô ($6 \times 5$ giá trị); chưa chạy trên đặc
        trưng hạt gạo của Phần D.
  \item \todo{P2, P3: hạn chế của phần mình (ví dụ: giới hạn thời gian ở C5, kích thước và độ đa dạng
        của dữ liệu gạo).}
\end{itemize}

\paragraph{Hướng phát triển.}
\begin{itemize}
  \item Thay cận $\nsv/n$ bằng các cận chặt hơn nhưng vẫn tính được sau một lần huấn luyện, như cận
        span hoặc cận bán kính--lề, để dùng trực tiếp cho việc chọn siêu tham số.
  \item Lặp lại khảo sát C4 trên đặc trưng hạt gạo để xem vùng $(C, \gamma)$ tốt có ổn định giữa các
        phiên chụp và thiết bị hay không.
  \item \todo{P2, P3: hướng phát triển cho phần cài đặt và ứng dụng.}
\end{itemize}
    % P1

\bibliographystyle{plainnat}
\bibliography{refs}

\appendix
% =====================================================================
%  Phụ lục: bảng phân công (P1)
% =====================================================================
\section{Phân công công việc}\label{sec:contrib}

\begin{table}[H]
\centering
\caption{Phân công công việc trong nhóm}\label{tab:contrib}
\begin{tabular}{@{}lll@{}}
\toprule
Thành viên & Vai trò & Phụ trách \\
\midrule
\todo{Họ tên P1} & Lý thuyết + báo cáo & B1--B6, C1--C4, tổng hợp báo cáo \\
\todo{Họ tên P2} & Cài đặt thuật toán & A1--A4, C5 \\
\todo{Họ tên P3} & Ứng dụng & D1--D5 \\
\bottomrule
\end{tabular}
\end{table}
  % P1
\end{document}
